\setcounter{ex}{0}

% --- CÂU 1 ---
\begin{ex}
	Cho phương trình $\cos\left(3x - \dfrac{\pi}{4}\right) = \dfrac{\sqrt{2}}{2}\ (*)$. Các mệnh đề sau đúng hay sai?
	\choiceTFt
	{\True Phương trình $(*)$ tương đương với $\cos\left(3x - \dfrac{\pi}{4}\right) = \cos\dfrac{\pi}{4}$}
	{\True Tập nghiệm của $(*)$ là $\left\{\dfrac{k2\pi}{3};\ \dfrac{\pi}{6} + \dfrac{k2\pi}{3},\ k \in \mathbb{Z}\right\}$}
	{\True Phương trình $(*)$ có đúng $4$ nghiệm thuộc đoạn $[0; \pi]$}
	{Nghiệm âm lớn nhất của $(*)$ bằng $-\dfrac{\pi}{6}$}
	\loigiai{
		\begin{itemchoice}
			\itemch \textbf{Đúng.} Vì $\cos\dfrac{\pi}{4} = \dfrac{\sqrt{2}}{2}$ nên phương trình $(*)$ tương đương với $\cos\left(3x - \dfrac{\pi}{4}\right) = \cos\dfrac{\pi}{4}$.
			\itemch \textbf{Đúng.} Ta có:
			$$\cos\left(3x - \dfrac{\pi}{4}\right) = \cos\dfrac{\pi}{4} \iff \hoac{3x - \dfrac{\pi}{4} &= \dfrac{\pi}{4} + k2\pi \\ 3x - \dfrac{\pi}{4} &= -\dfrac{\pi}{4} + k2\pi} \iff \hoac{x &= \dfrac{\pi}{6} + \dfrac{k2\pi}{3} \\ x &= \dfrac{k2\pi}{3}}\ (k \in \mathbb{Z}).$$
			Vậy tập nghiệm của $(*)$ là $\left\{\dfrac{k2\pi}{3};\ \dfrac{\pi}{6} + \dfrac{k2\pi}{3},\ k \in \mathbb{Z}\right\}$.
			\itemch \textbf{Đúng.} Trên đoạn $[0; \pi]$:
			\begin{itemize}
				\item Họ $x = \dfrac{k2\pi}{3} \in [0; \pi] \implies k \in \{0; 1\}$ (nghiệm $0; \dfrac{2\pi}{3}$).
				\item Họ $x = \dfrac{\pi}{6} + \dfrac{k2\pi}{3} \in [0; \pi] \implies k \in \{0; 1\}$ (nghiệm $\dfrac{\pi}{6}; \dfrac{5\pi}{6}$).
			\end{itemize}
			Vậy phương trình $(*)$ có đúng $4$ nghiệm thuộc đoạn $[0; \pi]$.
			\itemch \textbf{Sai.} Xét nghiệm âm ứng với $k = -1$:
			\begin{itemize}
				\item Họ $x = \dfrac{k2\pi}{3}$ cho nghiệm âm lớn nhất là $x = -\dfrac{2\pi}{3}$.
				\item Họ $x = \dfrac{\pi}{6} + \dfrac{k2\pi}{3}$ cho nghiệm âm lớn nhất là $x = \dfrac{\pi}{6} - \dfrac{2\pi}{3} = -\dfrac{\pi}{2}$.
			\end{itemize}
			Vậy nghiệm âm lớn nhất của $(*)$ bằng $-\dfrac{\pi}{2} \ne -\dfrac{\pi}{6}$.
		\end{itemchoice}
	}
\end{ex}

% --- CÂU 2 ---
\begin{ex}
	Cho hàm số $y = \dfrac{\cos x}{1 + \sin x}$. Các mệnh đề sau đúng hay sai?
	\choiceTFt
	{\True Tập xác định của hàm số là $\mathbb{R} \setminus \left\{-\dfrac{\pi}{2} + k2\pi,\ k \in \mathbb{Z}\right\}$}
	{\True Hàm số tuần hoàn với chu kì $2\pi$}
	{Hàm số đã cho là hàm số lẻ}
	{$y\left(\dfrac{\pi}{2}\right) = 1$}
	\loigiai{
		\begin{itemchoice}
			\itemch \textbf{Đúng.} Điều kiện xác định: $1 + \sin x \ne 0 \iff \sin x \ne -1 \iff x \ne -\dfrac{\pi}{2} + k2\pi\ (k \in \mathbb{Z})$.\\
			Vậy tập xác định của hàm số là $\mathscr{D} = \mathbb{R} \setminus \left\{-\dfrac{\pi}{2} + k2\pi,\ k \in \mathbb{Z}\right\}$.
			\itemch \textbf{Đúng.} Với mọi $x \in \mathscr{D}$, ta có $x + 2\pi \in \mathscr{D}$ và $y(x + 2\pi) = \dfrac{\cos(x + 2\pi)}{1 + \sin(x + 2\pi)} = \dfrac{\cos x}{1 + \sin x} = y(x)$, chu kì tuần hoàn là $2\pi$.
			\itemch \textbf{Sai.} Tập xác định $\mathscr{D}$ không phải là tập đối xứng (vì $\dfrac{\pi}{2} \in \mathscr{D}$ nhưng $-\dfrac{\pi}{2} \notin \mathscr{D}$) nên hàm số không phải là hàm số lẻ.
			\itemch \textbf{Sai.} Ta có $y\left(\dfrac{\pi}{2}\right) = \dfrac{\cos\frac{\pi}{2}}{1 + \sin\frac{\pi}{2}} = \dfrac{0}{1 + 1} = 0 \ne 1$.
		\end{itemchoice}
	}
\end{ex}

% --- CÂU 3 ---
\begin{ex}
	Cho cấp số cộng $(u_n)$ thỏa mãn $\heva{u_3 &= 11 \\ u_7 &= 27.}$ Các mệnh đề sau đúng hay sai?
	\choiceTFt
	{\True Công sai $d = 4$}
	{\True Số hạng đầu $u_1 = 3$}
	{Số $99$ là số hạng thứ $26$}
	{Tổng $15$ số hạng đầu $S_{15} = 420$}
	\loigiai{
		\begin{itemchoice}
			\itemch \textbf{Đúng.} Ta có hệ phương trình:
			$$\heva{u_3 &= u_1 + 2d = 11 \\ u_7 &= u_1 + 6d = 27} \iff \heva{u_1 &= 3 \\ d &= 4.}$$
			Vậy công sai $d = 4$.
			\itemch \textbf{Đúng.} Số hạng đầu $u_1 = 3$.
			\itemch \textbf{Sai.} Xét $u_n = u_1 + (n - 1)d = 99 \iff 3 + 4(n - 1) = 99 \iff n - 1 = 24 \iff n = 25 \ne 26$.
			\itemch \textbf{Sai.} Tổng $15$ số hạng đầu là:
			$$S_{15} = 15u_1 + \dfrac{15 \cdot 14}{2}d = 15 \cdot 3 + 105 \cdot 4 = 45 + 420 = 465 \ne 420.$$
		\end{itemchoice}
	}
\end{ex}

% --- CÂU 4 ---
\begin{ex}
	Cho hình chóp $S.ABCD$ có đáy $ABCD$ là hình bình hành tâm $O$, $M$ là trung điểm của $SC$, $F$ là giao điểm của $SD$ và mặt phẳng $(ABM)$ (hình vẽ). Các mệnh đề sau đúng hay sai?
	\begin{center}
		\begin{tikzpicture}[scale=0.85, font=\footnotesize, line join=round, line cap=round]
			\coordinate (A) at (0,0);
			\coordinate (B) at (3.2,0);
			\coordinate (D) at (1.3,1.1);
			\coordinate (C) at (4.5,1.1);
			\coordinate (O) at ($(A)!0.5!(C)$);
			\coordinate (S) at (2.2,3.4);
			\coordinate (M) at ($(S)!0.5!(C)$);
			\coordinate (F) at ($(S)!0.5!(D)$);
			\draw[thick] (S)--(A) (S)--(B) (S)--(C) (A)--(B)--(C) (B)--(M);
			\draw[dashed] (S)--(D) (A)--(D)--(C) (A)--(C) (B)--(D) (A)--(F)--(M);
			\foreach \p/\pos in {S/above, A/below left, B/below, C/right, D/left, O/below, M/above right, F/above left}
				\node[\pos] at (\p) {$\p$};
			\foreach \p in {S,A,B,C,D,O,M,F}
				\fill (\p) circle (1.2pt);
		\end{tikzpicture}
	\end{center}
	\choiceTFt
	{\True $F$ là trung điểm của $SD$}
	{\True $OM \parallel (SAB)$}
	{Thiết diện của hình chóp cắt bởi $(ABM)$ là hình bình hành}
	{\True Giao tuyến của $(SAD)$ và $(SBC)$ là đường thẳng đi qua $S$ và song song với $AD$}
	\loigiai{
		\begin{itemchoice}
			\itemch \textbf{Đúng.} Hai mặt phẳng $(ABM)$ và $(SCD)$ có điểm chung $M$ và chứa hai đường thẳng song song $AB \parallel CD$ nên giao tuyến của chúng là đường thẳng qua $M$ song song với $CD$, cắt $SD$ tại trung điểm $F$ của $SD$.
			\itemch \textbf{Đúng.} $OM$ là đường trung bình của tam giác $SAC \implies OM \parallel SA \subset (SAB) \implies OM \parallel (SAB)$.
			\itemch \textbf{Sai.} Thiết diện của hình chóp cắt bởi $(ABM)$ là tứ giác $ABMF$ có $MF \parallel AB$ và $MF = \dfrac{1}{2}CD = \dfrac{1}{2}AB < AB$ nên $ABMF$ là hình thang (không phải hình bình hành).
			\itemch \textbf{Đúng.} Hai mặt phẳng $(SAD)$ và $(SBC)$ có điểm chung $S$ và chứa hai đường thẳng song song $AD \parallel BC$ nên giao tuyến của chúng là đường thẳng đi qua $S$ và song song với $AD$.
		\end{itemchoice}
	}
\end{ex}
